%% fig-idea.tex — Slide 4 «Ідея» ★ the SYMBOL of the deck.
%% «Session as Node»: a peer appears on wg0 ONLY after passing 2FA.
%%
%% Single-path left→right story (simplified per user feedback R3):
%%   Клієнт+ключ →(недостатньо)→ 🔒 2FA → peer на wg0   [ONE state only]
%% Removal shown as a footnote only — no competing second peer node.
%% Assumes styles.tex pre-loaded.
\begin{tikzpicture}[x=1cm, y=1cm]
%%
%% ── SERVER NODE (left anchor of bus) ─────────────────────────────────────────
\node[wgboxsm, minimum width=3.0cm] (server) at (1.5, 0)
  {Сервер\\10.10.0.1};
%%
%% ── wg0 BUS — starts at server.east (no left-stub artefact) ─────────────────
\draw[wgbus] (server.east) -- (11.0, 0);
\node[font=\footnotesize\ttfamily, text=clGrey, above right] at (11.0, 0.1)
  {\texttt{wg0}};
%%
%% ── 2FA LOCK GATE ────────────────────────────────────────────────────────────
\node[lock, minimum width=2.0cm, minimum height=2.2cm] (lockbox) at (4.8, 0)
  {\textbf{2FA}};
%% Padlock shackle
\draw[clAccent, line width=1.8pt]
  ($(lockbox.north)+(+0.38,0)$)
  -- ++(0, 0.42)
  arc[radius=0.38, start angle=0, end angle=180]
  -- ++(0, -0.42);
%%
%% ── CLIENT NODE (left of lock, off-bus) ──────────────────────────────────────
\node[wgboxsm, minimum width=3.0cm] (client) at (2.8, -2.2)
  {Клієнт\\+ публічний ключ};
%% Dashed arrow client→lock (darkened locally for projector visibility)
\draw[->, >=Stealth, line width=1.4pt, color=black!62, dashed]
  (client.north) -- ($(lockbox.south west)!0.4!(lockbox.south east)$);
%% «ключа недостатньо» — centred between lock-gate bottom (y=-1.1) and
%% client-box top (y=-1.775); right edge ≈4.7 < 5, stays left of peer zone.
\node[font=\footnotesize\sffamily, text=black!62, align=center]
  at (3.0, -1.44) {ключа недостатньо};
%%
%% ── SESSION PEER: ONE state — solid, filled, on bus after the gate ───────────
\node[wgpeer, minimum width=3.4cm] (peer) at (8.0, 0)
  {10.10.0.x};
%%
%% ── ADD ACTION: solid accent arrow ↓ — session created ──────────────────────
\draw[wgadd] (8.0, 2.8) -- (peer.north);
\node[font=\small\ttfamily, text=clAccent, align=center]
  at (8.0, 3.15) {\texttt{wg set wg0 peer \ldots}};
%%
%% ── REMOVAL FOOTNOTE (small grey text only — no node box, no arrow) ──────────
%% Communicates the remove path without a competing second peer node.
%% Removal footnote moved to y=-2.3 — clearly BELOW the «ключа недостатньо» line.
\node[font=\footnotesize\sffamily, text=black!55, align=center]
  at (8.0, -2.3)
  {При завершенні сесії\\сервер прибирає користувача.};
%%
%% ── ANCHOR ────────────────────────────────────────────────────────────────────
\node[font=\large\sffamily\bfseries, text=clAccent, align=center]
  at (5.5, -3.5)
  {Без зміни протоколу.\quad Без PSK.};
%%
\end{tikzpicture}
